\section{Conclusion}

\label{sec:conclusion}

The Field protocol extends the metastable consensus family to multi-valued conflicts
through an electromagnetic-inspired framework. The field potential function provides
a rigorous Lyapunov proof of convergence, and the phase transition analysis characterizes
the parameter regimes where consensus is achievable. The stake-weighted extension via
field strength gradients naturally integrates proof-of-stake economics with the consensus
dynamics. Field provides the theoretical foundation for multi-valued decision-making
in the Lux network, including validator set elections, governance votes, and multi-party
auction resolution.

\begin{thebibliography}{99}
\bibitem{rocket2019wave} T. Rocket et al., ``Scalable and probabilistic leaderless BFT consensus through metastability,'' 2019.
\bibitem{kelling2020wave} Z. Kelling, ``Wave protocol: Decision stability and confidence monotonicity,'' Lux Industries, Inc., 2020.
\bibitem{kelling2021flare} Z. Kelling, ``Flare protocol: Certificate skip exclusivity with honest support,'' Lux Industries, Inc., 2021.
\bibitem{ising1925beitrag} E. Ising, ``Beitrag zur Theorie des Ferromagnetismus,'' \emph{Zeitschrift f\"ur Physik}, 1925.
\bibitem{potts1952some} R. Potts, ``Some generalized order-disorder transformations,'' \emph{Mathematical Proceedings of the Cambridge Philosophical Society}, 1952.
\bibitem{glauber1963time} R. Glauber, ``Time-dependent statistics of the Ising model,'' \emph{J. Math. Phys.}, 1963.
\bibitem{landau1980statistical} L. Landau and E. Lifshitz, \emph{Statistical Physics}, 3rd ed., Pergamon Press, 1980.
\end{thebibliography}

